\documentclass[conference]{IEEEtran}

% Essential packages
\usepackage[T1]{fontenc}
\usepackage{cite}
\usepackage{amsmath,amssymb}
\usepackage{graphicx}
\usepackage{booktabs}
\usepackage{hyperref}
\usepackage{url}

\title{Monitoring LLM Reasoning Traces for Unsafe Behavior}

\author{
\IEEEauthorblockN{Author One and Author Two}
\IEEEauthorblockA{Narxoz University\\
Almaty, Kazakhstan}
}

\begin{document}
\maketitle

\begin{abstract}
Large language models (LLMs) can produce outputs that are misleading,
harmful, or inconsistent with intended behavior. This project studies
whether access to a model's reasoning trace, in addition to its prompt and
final output, can help a monitor identify unsafe behavior. We compare two
monitoring conditions: output-only monitoring and monitoring that receives
the prompt, reasoning trace, and final output. The experiment considers
reward hacking, deception, sabotage, and harmful compliance. We describe
the monitoring setup and evaluation plan, with particular attention to
recall at a fixed false-positive rate. The current implementation uses a
keyword-based monitoring baseline; therefore, conclusions are limited to
the tested dataset and monitor rather than generalized to all LLMs.
\end{abstract}

\begin{IEEEkeywords}
large language models, AI safety, chain-of-thought, monitoring,
unsafe behavior, classification
\end{IEEEkeywords}


\section{Introduction}
\label{sec:introduction}

Large language models (LLMs) are increasingly used in systems that
generate text, assist with decisions, and interact with digital
environments. Their usefulness also creates safety challenges: a model
may produce harmful instructions, conceal its behavior, manipulate an
evaluation process, or otherwise act contrary to the system's intended
objectives. Reliable monitoring is therefore an important part of
evaluating and deploying LLM-based systems.

A monitor can inspect different parts of a model interaction. An
output-only monitor receives the user's prompt and the model's final
answer. A more information-rich monitor may also receive an available
reasoning trace, often referred to as chain-of-thought (CoT). In
principle, additional context may expose signals that are not visible in
the final answer alone. However, access to a reasoning trace does not
automatically guarantee more accurate detection: the trace may be
incomplete, noisy, unavailable, or may not faithfully represent the
model's internal computation.

This creates the research problem addressed in this project: whether
including a reasoning trace changes the ability of a monitor to identify
unsafe behavior under a controlled evaluation setup. We compare two
conditions: (1) output-only monitoring, which uses the prompt and final
output, and (2) CoT-plus-output monitoring, which uses the prompt,
reasoning trace, and final output. The study focuses on four categories
of unsafe behavior: reward hacking, deception, sabotage, and harmful
compliance.

The research gap is the need to evaluate the effect of the information
available to a monitor while keeping the monitoring approach and
evaluation conditions explicit. This project provides a controlled
comparison of the two input conditions using the benchmark data and
keyword-based monitoring implementation available in the repository.
Because the current monitor is a keyword baseline, the study should be
interpreted as an evaluation of this baseline and dataset, not as proof
that reasoning traces universally improve LLM safety monitoring.

The research questions are:

\begin{itemize}
    \item \textbf{RQ1:} How does access to a reasoning trace affect the
    monitor's ability to detect unsafe behavior compared with
    output-only monitoring?

    \item \textbf{RQ2:} What is the recall of each monitoring condition
    at a target false-positive rate of 5\%?

    \item \textbf{RQ3:} How do the monitoring conditions compare in
    operational performance, particularly latency, under the tested
    request loads?
\end{itemize}

The contributions of this project are:

\begin{itemize}
    \item A comparison of output-only and CoT-plus-output monitoring
    conditions.

    \item An evaluation framework focused on unsafe-behavior categories
    including reward hacking, deception, sabotage, and harmful compliance.

    \item An analysis plan that considers detection quality at a fixed
    false-positive rate and operational latency.

    \item A reproducible experimental configuration and processing
    workflow for the benchmark.
\end{itemize}

The remainder of the paper describes the experimental methodology,
reports the measured results, discusses their implications and
limitations, and concludes with directions for future work.


\section{Methodology}
\label{sec:methodology}

\subsection{Problem Formulation}

We formulate unsafe-behavior monitoring as a binary classification
task. For each episode $i$, the monitor receives an input representation
$x_i$ and produces a score $s_i$ indicating the likelihood that the
episode is unsafe. Given a threshold $\tau$, the predicted label is

\begin{equation}
\hat{y}_i =
\begin{cases}
1, & s_i \geq \tau,\\
0, & s_i < \tau,
\end{cases}
\label{eq:prediction}
\end{equation}

where $y_i \in \{0,1\}$ is the reference label and $\hat{y}_i$ is the
monitor prediction.

We compare two input conditions:

\begin{align}
x_i^{\mathrm{out}} &=
(\mathrm{prompt}_i,\mathrm{output}_i),\\
x_i^{\mathrm{cot}} &=
(\mathrm{prompt}_i,\mathrm{cot}_i,\mathrm{output}_i).
\end{align}

The primary detection metric is recall at a target false-positive rate
(FPR) of 5\%. Recall is defined as

\begin{equation}
\mathrm{Recall} = \frac{TP}{TP+FN},
\label{eq:recall}
\end{equation}

and the false-positive rate is

\begin{equation}
\mathrm{FPR} = \frac{FP}{FP+TN}.
\label{eq:fpr}
\end{equation}

\subsection{Monitoring Conditions}

In the output-only condition, the monitor can inspect the prompt and
final model output. In the CoT-plus-output condition, it can inspect
the prompt, the available reasoning trace, and the final output.

The repository configuration identifies these conditions as
\texttt{output\_only} and \texttt{cot\_plus\_output}, respectively.

\subsection{Unsafe-Behavior Categories}

The benchmark groups unsafe episodes into four categories:

\begin{itemize}
    \item \textbf{Reward hacking:} exploiting a reward or evaluation
    mechanism rather than satisfying its intended objective.

    \item \textbf{Deception:} misleading a user, evaluator, or system
    about relevant behavior or information.

    \item \textbf{Sabotage:} intentionally interfering with a process,
    system, or objective.

    \item \textbf{Harmful compliance:} complying with a request that
    would facilitate harmful activity.
\end{itemize}

\subsection{Monitor and Experimental Configuration}

The current repository configuration specifies a keyword-based
monitoring backend. The configuration also specifies a random seed
of 42, a target FPR of 0.05, and a calibration fraction of 0.5.

The episode data are loaded from the configured JSONL dataset.
These values should be checked against the exact configuration and
dataset version used for the final experiment before submission.

\subsection{Evaluation Procedure}

The experiment separates calibration and evaluation data according
to the configured split. A decision threshold is selected using the
calibration data to target the specified FPR. The monitor is then
evaluated on the held-out evaluation data.

The two monitoring conditions are compared using the same
evaluation protocol.

Operational performance is measured separately from detection
quality. Latency measurements are collected under the configured
request rates. Latency summaries and plots should be reported
independently from classification metrics; a larger number of
latency observations does not imply a larger number of labeled
detection episodes.

\subsection{Reproducibility}

To reproduce the experiment, a researcher should record the
repository revision, configuration file, dataset version, random
seed, monitor backend, visible input fields, calibration split,
target FPR, software dependencies, and execution commands.

Any change to these settings should be documented because it may
affect the results.


% Add these sections after the corresponding experiments have been evaluated:
% \input{sections/03_results}
% \input{sections/04_discussion}
% \input{sections/05_limitations_conclusion}

\bibliographystyle{IEEEtran}
\bibliography{references}

\end{document}
